\documentclass{article}
\usepackage{amsmath,amssymb,amsthm}
\usepackage[margin=1in]{geometry}

\newtheorem{theorem}{Theorem}
\newtheorem{lemma}[theorem]{Lemma}

\DeclareMathOperator{\lcm}{lcm}

\begin{document}

\title{IMO 2026 Problem 1 --- GCD/LCM Blackboard}
\author{}
\date{}
\maketitle

\section*{Problem Statement}

There are 2026 integers greater than 1 written on a blackboard, not necessarily different.
In a move, Confucius chooses two integers $m>1$ and $n>1$ from different places on the blackboard
and replaces these two integers with
\[
\gcd(m,n) \quad\text{and}\quad \frac{\lcm(m,n)}{\gcd(m,n)}.
\]
He continues to make moves while it is possible to do so.

\begin{enumerate}
\item[(a)] Prove that, regardless of the choices of Confucius, after finitely many moves,
           exactly one integer $M$ on the blackboard is greater than 1.
\item[(b)] Prove that the value of $M$ does not depend on the choices of Confucius.
\end{enumerate}

\section{Notation and Preliminaries}

For a positive integer $n$ and a prime $p$, let $v_p(n)$ denote the $p$-adic valuation:
the unique nonnegative integer such that $p^{v_p(n)}\mid n$ and $p^{v_p(n)+1}\nmid n$;
set $v_p(1)=0$. By the fundamental theorem of arithmetic, every positive integer $n$ is
uniquely $n=\prod_p p^{v_p(n)}$ (finitely many $p$ with $v_p(n)>0$); consequently two
positive integers are equal iff their $p$-adic valuations agree for every prime $p$.

For a finite multiset $S$ of integers, define $\gcd(S)$ as the iterated gcd of its elements,
with the conventions $\gcd(\emptyset)=0$, $\gcd(0,0)=0$, and $\gcd(x,0)=x$ for $x\geq 0$;
the value is independent of the iteration order by associativity and commutativity of $\gcd$.

\section{Lemmas}

\begin{lemma}[Exponent formulas for gcd and lcm]\label{lem:vpgcd}
For a prime $p$ and positive integers $a,b$,
\[
v_p(\gcd(a,b))=\min(v_p(a),v_p(b)),\qquad
v_p(\lcm(a,b))=\max(v_p(a),v_p(b)).
\]
\end{lemma}
\begin{proof}
Write $a=\prod_q q^{\alpha_q}$, $b=\prod_q q^{\beta_q}$ with $\alpha_q=v_q(a)$,
$\beta_q=v_q(b)\geq 0$. By unique factorization, $d\mid a$ and $d\mid b$ iff
$v_q(d)\leq\min(\alpha_q,\beta_q)$ for every prime $q$; hence the greatest common divisor of
$a,b$ has exponents $\delta_q=\min(\alpha_q,\beta_q)$, and, similarly, the least common
multiple has exponents $\max(\alpha_q,\beta_q)$. Thus
$v_p(\gcd(a,b))=\min(v_p(a),v_p(b))$ and $v_p(\lcm(a,b))=\max(v_p(a),v_p(b))$ for each prime $p$.
\end{proof}

\begin{lemma}[Identity $\lcm\cdot\gcd=mn$]\label{lem:lcmgcd}
For positive integers $m,n$ with $g=\gcd(m,n)$,
\[
\lcm(m,n)\cdot\gcd(m,n)=mn,\qquad\text{so}\qquad \frac{\lcm(m,n)}{\gcd(m,n)}=\frac{m}{g}\cdot\frac{n}{g}\in\mathbb{Z}_{>0}.
\]
Under a legal move, the product $P$ of all board entries changes from $P$ to $P/g$,
and the new pair $(g,\lcm(m,n)/g)$ always contains at least one entry $>1$.
\end{lemma}
\begin{proof}
For each prime $p$, Lemma~\ref{lem:vpgcd} gives
$v_p(\lcm(m,n)\gcd(m,n))=\max(v_p(m),v_p(n))+\min(v_p(m),v_p(n))=v_p(mn)$, so the two positive
integers are equal by unique factorization. Hence
$\lcm(m,n)/\gcd(m,n)=mn/g^2=(m/g)(n/g)$, a positive integer because $g\mid m$ and $g\mid n$.
The pair product is $g\cdot(mn/g^2)=mn/g$, so the total product $P$ is divided by $g$.
If $g>1$, the first new entry $g>1$; if $g=1$, the new pair is $(1,mn)$ with $mn\geq 4$
(since $m,n\geq 2$). In both cases the new pair contains an entry $>1$.
\end{proof}

\begin{lemma}[Per-prime exponent dynamics of a move]\label{lem:expdyn}
Fix a prime $p$. A legal move on entries $m,n$ with $(x,y)=(v_p(m),v_p(n))$ replaces
the exponent pair by $(\min(x,y),|x-y|)$. The $p$-valuations of all other entries are
unchanged, and the same rule holds uniformly for every prime.
\end{lemma}
\begin{proof}
By Lemma~\ref{lem:lcmgcd}, $h=\lcm(m,n)/\gcd(m,n)$ is a positive integer, so $v_p(h)$ is well
defined; writing $g=\gcd(m,n)$, Lemma~\ref{lem:vpgcd} gives
\[
v_p(h)=v_p(\lcm(m,n))-v_p(g)=\max(x,y)-\min(x,y)=|x-y|.
\]
The first new entry is $g$, with $v_p(g)=\min(x,y)$, and entries not moved keep their
valuations. The argument uses only that $p$ is prime, so it applies uniformly to every
prime (for a prime dividing neither $m$ nor $n$, the pair $(0,0)$ is fixed).
\end{proof}

\begin{lemma}[Euclidean step identity for gcd]\label{lem:eucid}
For all nonnegative integers $x,y$, with $\gcd(0,0)=0$ and $\gcd(x,0)=x$,
\[
\gcd(x,y)=\gcd(\min(x,y),\,|x-y|).
\]
\end{lemma}
\begin{proof}
Assume $x\leq y$; the other case is symmetric. Then $\min(x,y)=x$ and $|x-y|=y-x$.
The ideals $(x,y)$ and $(x,y-x)$ are equal: $y=(y-x)+x\in(x,y-x)$, and $y-x\in(x,y)$.
In $\mathbb{Z}$ every ideal is principal, and the unique nonnegative generator of $(a,b)$
is $\gcd(a,b)$, so $\gcd(x,y)=\gcd(x,y-x)$. The conventions $\gcd(0,0)=0$ and
$\gcd(x,0)=x$ follow from $(0,0)=0\mathbb{Z}$ and $(x,0)=x\mathbb{Z}$.
\end{proof}

\begin{lemma}[Multiset gcd preservation]\label{lem:mgcdpres}
Let $S$ be a finite multiset of integers with $|S|\geq 1$ and let $S'$ be obtained from $S$
by replacing two elements $a,b$ at distinct positions by $c,d$ such that
$\gcd(c,d)=\gcd(a,b)$. Then $\gcd(S')=\gcd(S)$.
\end{lemma}
\begin{proof}
Let $R$ be the multiset obtained from $S$ by deleting the two positions holding $a$ and $b$.
By the grouping identity for the iterated gcd,
$\gcd(S)=\gcd(\gcd(a,b),\gcd(R))$ and $\gcd(S')=\gcd(\gcd(c,d),\gcd(R))$, which are equal by
hypothesis. If $|S|=2$ then $R=\emptyset$, and the convention $\gcd(\emptyset)=0$ together
with $\gcd(x,0)=x$ covers this case.
\end{proof}

\begin{lemma}[Exponent-gcd invariant --- part (b)]\label{lem:inv}
For every prime $p$, the quantity
\[
G_p \;:=\; \gcd\{\,v_p(a)\mid a\text{ an entry of the board}\,\}
\]
is invariant under every legal move, with the convention $\gcd(0,0)=0$ on exponents.
\end{lemma}
\begin{proof}
At every board state all entries are positive integers, so
$M_p(B)=\{v_p(a):a\text{ on board}\}$ is a finite multiset of nonnegative integers and its
iterated gcd $G_p(B)$ is well defined with the stated conventions.

Consider a legal move on entries $m,n$ with $(x,y)=(v_p(m),v_p(n))$. By
Lemma~\ref{lem:expdyn}, the pair $(x,y)$ is replaced by $(\min(x,y),|x-y|)$ and all other
$p$-valuations are unchanged, so
\[
M_p(B') = (M_p(B)\setminus\{x,y\})\cup\{\min(x,y),|x-y|\}.
\]
By Lemma~\ref{lem:eucid}, $\gcd(\min(x,y),|x-y|)=\gcd(x,y)$, so the replacement pair has the
same gcd as the pair it replaces; applying Lemma~\ref{lem:mgcdpres} to $M_p(B)$ with
$a=x$, $b=y$, $c=\min(x,y)$, $d=|x-y|$ gives $G_p(B')=G_p(B)$.

Thus $G_p$ is unchanged by every move; inductively over the move sequence, $G_p$ is constant
throughout any play and equals its initial value.
\end{proof}

\begin{lemma}[Termination --- part (a)]\label{lem:term}
For any board of $k\geq 1$ integers~${}>1$ (with $k=2026$ in the problem) and any legal play:
\begin{enumerate}
\item[(i)] The number $C$ of entries $>1$ satisfies $C\geq 1$ throughout and never increases.
\item[(ii)] The pair $(P,C)$, where $P$ is the product of all entries, strictly decreases
            in lexicographic order after every move.
\item[(iii)] Every play has at most $(k-1)+\lfloor\log_2 P_0\rfloor$ moves and is therefore finite.
\item[(iv)] The terminal board has exactly one entry $>1$.
\end{enumerate}
\end{lemma}
\begin{proof}
For a move on entries $m,n$ with $g=\gcd(m,n)$, Lemma~\ref{lem:lcmgcd} gives that the new pair
$(g,\lcm(m,n)/g)$ has product $mn/g$ and contains an entry $>1$.

\textbf{Case $g>1$:} the total product is divided by $g\geq 2$, so $P$ strictly decreases, and
$C$ does not increase (the new pair contributes at most two entries $>1$).

\textbf{Case $g=1$:} $P$ is unchanged, and the new pair is $(1,mn)$ with $mn\geq 4>1$, so $C$
decreases by exactly $1$.

Hence $C$ never increases; and $C\geq 1$ is preserved, because every new pair contains an entry
$>1$ (initially $C=k\geq 1$). Therefore $(P,C)$ strictly decreases in lexicographic order after
every move, proving (i) and (ii).

Moves with $g>1$ divide $P$ by at least $2$ each, and $P\geq 1$ throughout, so there are at most
$\lfloor\log_2 P_0\rfloor$ of them. Moves with $g=1$ decrease $C$ by exactly $1$, and $C\geq 1$,
so there are at most $k-1$ of them. Every play thus has at most $(k-1)+\lfloor\log_2 P_0\rfloor$
moves and is finite, proving (iii).

A move is possible exactly when $C\geq 2$ (two entries $>1$ at distinct places); hence a maximal
play stops only at $C=1$, i.e.\ the terminal board has exactly one entry $>1$, proving (iv).
\end{proof}

\begin{lemma}[Terminal valuation equality]\label{lem:termval}
At the terminal board of any legal play, with $M$ the unique entry $>1$,
for every prime $p$:
\[
v_p(M)=G_p,
\]
where $G_p$ is the invariant from Lemma~\ref{lem:inv}, equal to the gcd of the
$p$-adic valuations of the \emph{initial} entries.
\end{lemma}
\begin{proof}
By Lemma~\ref{lem:term}(iv), the terminal board consists of $M$ together with $2025$ entries
equal to $1$. The exponent multiset at $p$ is therefore $\{v_p(M),0,\dots,0\}$ ($2025$ zeros),
since $v_p(1)=0$, and its gcd is $v_p(M)$ by repeated use of $\gcd(x,0)=x$ (the convention
$\gcd(0,0)=0$ covers the case $v_p(M)=0$). By Lemma~\ref{lem:inv}, this terminal value equals
the initial value of $G_p$, i.e.\ the gcd of the $p$-adic valuations of the initial $2026$
entries.
\end{proof}

\section{Proof of the Main Theorem}

\begin{theorem}[IMO 2026 Problem 1]
For the initial board of $2026$ integers ${}>1$ and the legal moves as defined:
\begin{enumerate}
\item[(a)] Every play terminates after finitely many moves, and the terminal board
           contains exactly one integer $M>1$.
\item[(b)] The value of $M$ is independent of the choices made during the play; explicitly,
           $v_p(M)=G_p$ for every prime $p$, where $G_p=\gcd_i v_p(a_i)$ is the gcd of the
           $p$-adic valuations of the initial entries, so
           \[
           M=\prod_{p} p^{\,G_p}
           \]
           (the ``exponent-gcd number'' of the initial multiset), a finite product
           determined solely by the initial board.
\end{enumerate}
\end{theorem}
\begin{proof}
\textbf{Part (a).} Instantiate Lemma~\ref{lem:term} with $k=2026$. All preconditions hold:
there are $2026$ entries all $>1$, a legal move selects two entries $>1$ at distinct places,
and play continues while possible. Hence every play is finite and its terminal board has
exactly one entry $M>1$.

\textbf{Part (b).} At the terminal board, the $2025$ entries other than $M$ are positive
integers not exceeding $1$, hence equal to $1$. By Lemma~\ref{lem:termval},
$v_p(M)=G_p$ for every prime $p$, where $G_p$ is the gcd of the $p$-adic valuations of the
initial entries. Unique prime factorization gives
$M=\prod_p p^{v_p(M)}=\prod_p p^{G_p}$, a finite product: $G_p>0$ only for primes dividing
some initial entry (all other $p$-valuations are $0$, so $G_p=\gcd(0,\dots,0)=0$ and
$p^{G_p}=1$). If two plays terminate at $M_1$ and $M_2$, then $v_p(M_1)=G_p=v_p(M_2)$ for all
$p$, so $M_1=M_2$ by unique factorization; hence $M$ is determined solely by the initial
multiset and does not depend on the choices of Confucius.

Finally, $M>1$: some initial entry is $\geq 2$ and has a prime divisor $p$ with
$v_p(\text{that entry})\geq 1$, so $G_p\geq 1$ and $p^{G_p}\mid M$, consistent with part (a).
\end{proof}

\section*{Remark}

The value $M=\prod_p p^{G_p}$ is in general \textbf{not} equal to the numeric gcd of the
initial entries. For example, on initial board $(4,8)$, one possible play is
$(4,8)\to(4,2)\to(2,2)\to(2,1)$, terminating with $M=2$, while $\gcd(4,8)=4$. The quantities
differ because the $p$-adic valuation of the numeric gcd is $\min_i v_p(a_i)$, whereas
$G_p=\gcd_i v_p(a_i)$ is the gcd of the valuations, not their minimum. The statement of part
(b) asks only for independence of choices, which the exponent-gcd invariant proves. The
theorem above gives the correct explicit formula.

\end{document}
